\documentclass[UTF8,11pt]{ctexart}
\usepackage[a4paper,margin=24mm]{geometry}
\usepackage{amsmath,amssymb,amsthm,mathtools}
\usepackage[all]{xy}
\usepackage{xurl,hyperref,enumitem}
\hypersetup{hidelinks,breaklinks=true}
\setlength{\emergencystretch}{3em}
\newtheorem*{theorem}{Theorem}
\newtheorem*{lemma}{Lemma}
\newtheorem*{proposition}{Proposition}
\newtheorem*{corollary}{Corollary}
\theoremstyle{definition}
\newtheorem*{definition}{Definition}
\newtheorem*{remark}{Remark}
\DeclareMathOperator{\Hom}{Hom}
\DeclareMathOperator{\Ext}{Ext}
\DeclareMathOperator{\Tor}{Tor}
\DeclareMathOperator{\Spec}{Spec}
\DeclareMathOperator{\supp}{supp}
\DeclareMathOperator{\im}{im}
\DeclareMathOperator{\id}{id}
\DeclareMathOperator{\Tr}{Tr}
\DeclareMathOperator{\rank}{rank}
\newcommand{\R}{\mathbb R}
\newcommand{\C}{\mathbb C}
\newcommand{\Z}{\mathbb Z}
\newcommand{\N}{\mathbb N}
\newcommand{\Q}{\mathbb Q}
\begin{document}
\section*{029\quad 正则局部环的有限整体维数刻画}
\subsection*{来源与收录范围}
The Stacks Project Authors, \emph{The Stacks Project}. Chapter 10, Proposition 10.110.5；收录 10.110.1、10.110.3--4 及正向消解构造所用 10.104.8--9、10.106.5--6。
\par\url{https://stacks.math.columbia.edu/tag/065U}
\par\url{https://stacks.math.columbia.edu/tag/00OC}
\par\url{https://stacks.math.columbia.edu/tag/00NT}
\par\url{https://stacks.math.columbia.edu/tag/00NS}
\par\url{https://stacks.math.columbia.edu/tag/00NG}
\par\url{https://stacks.math.columbia.edu/tag/00NE}
\par 使用 2026-09-28 实际访问的在线正文；未取得网页构建的固定版本号。以网页数学 TeX 整理，不把本摘录称为整章下载原件。
\subsection*{作者命题与人类证明}
\begin{lemma}[10.106.5]
Let $R$ be a Noetherian local ring. Let $x\in\mathfrak m$. Let $M$ be a finite $R$-module such that $x$ is a nonzerodivisor on $M$ and $M/xM$ is free over $R/xR$. Then $M$ is free over $R$.
\end{lemma}
\begin{proof}
Let $m_1,\ldots,m_r$ be elements of $M$ which map to a $R/xR$-basis of $M/xM$. By Nakayama's Lemma 10.20.1 $m_1,\ldots,m_r$ generate $M$. If $\sum a_im_i=0$ is a relation, then $a_i\in xR$ for all $i$. Hence $a_i=b_ix$ for some $b_i\in R$. Hence the kernel $K$ of $R^r\to M$ satisfies $xK=K$ and hence is zero by Nakayama's lemma.
\end{proof}
\begin{lemma}[10.106.6]
Let $R$ be a regular local ring. Any maximal Cohen-Macaulay module over $R$ is free.
\end{lemma}
\begin{proof}
Let $M$ be a maximal Cohen-Macaulay module over $R$. Let $x\in\mathfrak m$ be part of a regular sequence generating $\mathfrak m$. Then $x$ is a nonzerodivisor on $M$ by Proposition 10.103.4, and $M/xM$ is a maximal Cohen-Macaulay module over $R/xR$. By induction on $\dim(R)$ we see that $M/xM$ is free. We win by Lemma 10.106.5.
\end{proof}
\begin{lemma}[10.104.8]
Let $R$ be a Noetherian local Cohen-Macaulay ring of dimension $d$. Let $0\to K\to R^{\oplus n}\to M\to0$ be an exact sequence of $R$-modules. Then either $M=0$, or $\operatorname{depth}(K)>\operatorname{depth}(M)$, or $\operatorname{depth}(K)=\operatorname{depth}(M)=d$.
\end{lemma}
\begin{proof}
This is a special case of Lemma 10.72.6.
\end{proof}
\begin{lemma}[10.104.9]
Let $R$ be a local Noetherian Cohen-Macaulay ring of dimension $d$. Let $M$ be a finite $R$-module of depth $e$. There exists an exact complex
\[
0\to K\to F_{d-e-1}\to\ldots\to F_0\to M\to0
\]
with each $F_i$ finite free and $K$ maximal Cohen-Macaulay.
\end{lemma}
\begin{proof}
Immediate from the definition and Lemma 10.104.8.
\end{proof}
\begin{proposition}[10.110.1]
Let $R$ be a regular local ring of dimension $d$. Every finite $R$-module $M$ of depth $e$ has a finite free resolution
\[
0\to F_{d-e}\to\ldots\to F_0\to M\to0.
\]
In particular a regular local ring has global dimension $\leq d$.
\end{proposition}
\begin{proof}
The first part holds in view of Lemma 10.106.6 and Lemma 10.104.9. The last part follows from this and Lemma 10.109.12.
\end{proof}
\begin{lemma}[10.110.3]
Suppose that $R$ is a Noetherian local ring with maximal ideal $\mathfrak m$ and residue field $\kappa$. In this case the projective dimension of $\kappa$ is $\geq\dim_\kappa\mathfrak m/\mathfrak m^2$.
\end{lemma}
\begin{proof}
Let $x_1,\ldots,x_n$ be elements of $\mathfrak m$ whose images in $\mathfrak m/\mathfrak m^2$ form a basis. Consider the Koszul complex on $x_1,\ldots,x_n$. This is the complex
\[
0\to\wedge^nR^n\to\wedge^{n-1}R^n\to\wedge^{n-2}R^n\to\ldots\to\wedge^iR^n\to\ldots\to R^n\to R
\]
with maps given by
\[
e_{j_1}\wedge\ldots\wedge e_{j_i}\longmapsto
\sum_{a=1}^i(-1)^{a+1}x_{j_a}e_{j_1}\wedge\ldots\wedge\widehat{e_{j_a}}\wedge\ldots\wedge e_{j_i}.
\]
It is easy to see that this is a complex $K_\bullet(R,x_\bullet)$. Note that the cokernel of the last map of $K_\bullet(R,x_\bullet)$ is $\kappa$ by Lemma 10.20.1 part (8).

If $\kappa$ has finite projective dimension $d$, then we can find a resolution $F_\bullet\to\kappa$ by finite free $R$-modules of length $d$ (Lemma 10.109.7). By Lemma 10.102.2 we may assume all the maps in the complex $F_\bullet$ have the property that $\im(F_i\to F_{i-1})\subset\mathfrak mF_{i-1}$, because removing a trivial summand from the resolution can at worst shorten the resolution. By Lemma 10.71.4 we can find a map of complexes $\alpha:K_\bullet(R,x_\bullet)\to F_\bullet$ inducing the identity on $\kappa$. We will prove by induction that the maps $\alpha_i:\wedge^iR^n=K_i(R,x_\bullet)\to F_i$ have the property that $\alpha_i\otimes\kappa:\wedge^i\kappa^n\to F_i\otimes\kappa$ are injective. This shows that $F_n\ne0$ and hence $d\geq n$ as desired.

The result is clear for $i=0$ because the composition $R\xrightarrow{\alpha_0}F_0\to\kappa$ is nonzero. Note that $F_0$ must have rank $1$ since otherwise the map $F_1\to F_0$ whose cokernel is a single copy of $\kappa$ cannot have image contained in $\mathfrak mF_0$.

Next we check the case $i=1$ as we feel that it is instructive; the reader can skip this as the induction step will deduce the $i=1$ case from the case $i=0$. We saw above that $F_0=R$ and $F_1\to F_0=R$ has image $\mathfrak m$. We have a commutative diagram
\[
\begin{matrix}
R^n&=&K_1(R,x_\bullet)&\to&K_0(R,x_\bullet)&=&R\\
&&\downarrow&&\downarrow&&\downarrow\\
&&F_1&\to&F_0&=&R
\end{matrix}
\]
where the rightmost vertical arrow is given by multiplication by a unit. Hence we see that the image of the composition $R^n\to F_1\to F_0=R$ is also equal to $\mathfrak m$. Thus the map $R^n\otimes\kappa\to F_1\otimes\kappa$ has to be injective since $\dim_\kappa(\mathfrak m/\mathfrak m^2)=n$.

Let $i\geq1$ and assume injectivity of $\alpha_j\otimes\kappa$ has been proved for all $j\leq i-1$. Consider the commutative diagram
\[
\begin{matrix}
\wedge^iR^n&=&K_i(R,x_\bullet)&\to&K_{i-1}(R,x_\bullet)&=&\wedge^{i-1}R^n\\
&&\downarrow&&\downarrow&&\\
&&F_i&\to&F_{i-1}&&
\end{matrix}
\]
We know that $\wedge^{i-1}\kappa^n\to F_{i-1}\otimes\kappa$ is injective. This proves that $\wedge^{i-1}\kappa^n\otimes_\kappa\mathfrak m/\mathfrak m^2\to F_{i-1}\otimes\mathfrak m/\mathfrak m^2$ is injective. Also, by our choice of the complex, $F_i$ maps into $\mathfrak mF_{i-1}$, and similarly for the Koszul complex. Hence we get a commutative diagram
\[
\begin{matrix}
\wedge^i\kappa^n&\to&\wedge^{i-1}\kappa^n\otimes\mathfrak m/\mathfrak m^2\\
\downarrow&&\downarrow\\
F_i\otimes\kappa&\to&F_{i-1}\otimes\mathfrak m/\mathfrak m^2
\end{matrix}
\]
At this point it suffices to verify the map $\wedge^i\kappa^n\to\wedge^{i-1}\kappa^n\otimes\mathfrak m/\mathfrak m^2$ is injective, which can be done by hand.
\end{proof}
\begin{lemma}[10.110.4]
Let $R$ be a Noetherian local ring. Suppose that the residue field $\kappa$ has finite projective dimension $n$ over $R$. In this case $\dim(R)\geq n$.
\end{lemma}
\begin{proof}
Let $F_\bullet$ be a finite resolution of $\kappa$ by finite free $R$-modules (Lemma 10.109.7). By Lemma 10.102.2 we may assume all the maps in the complex $F_\bullet$ have to property that $\im(F_i\to F_{i-1})\subset\mathfrak mF_{i-1}$, because removing a trivial summand from the resolution can at worst shorten the resolution. Say $F_n\ne0$ and $F_i=0$ for $i>n$, so that the projective dimension of $\kappa$ is $n$. By Proposition 10.102.9 we see that $\operatorname{depth}_{I(\varphi_n)}(R)\geq n$ since $I(\varphi_n)$ cannot equal $R$ by our choice of the complex. Thus by Lemma 10.72.3 also $\dim(R)\geq n$.
\end{proof}
\begin{proposition}[10.110.5]
Let $(R,\mathfrak m,\kappa)$ be a Noetherian local ring. The following are equivalent:
\begin{enumerate}
\item $\kappa$ has finite projective dimension as an $R$-module;
\item $R$ has finite global dimension;
\item $R$ is a regular local ring.
\end{enumerate}
Moreover, in this case the global dimension of $R$ equals $\dim(R)=\dim_\kappa(\mathfrak m/\mathfrak m^2)$.
\end{proposition}
\begin{proof}
We have (3) $\Rightarrow$ (2) by Proposition 10.110.1. The implication (2) $\Rightarrow$ (1) is trivial. Assume (1). By Lemmas 10.110.3 and 10.110.4 we see that $\dim(R)\geq\dim_\kappa(\mathfrak m/\mathfrak m^2)$. Thus $R$ is regular, see Definition 10.60.10 and the discussion preceding it. Assume the equivalent conditions (1)--(3) hold. By Proposition 10.110.1 the global dimension of $R$ is at most $\dim(R)$ and by Lemma 10.110.3 it is at least $\dim_\kappa(\mathfrak m/\mathfrak m^2)$. Thus the stated equality holds.
\end{proof}

\subsection*{整理说明（非作者正文）}
主目标为正则局部环、剩余域有限投射维数和有限整体维数的等价及维数相等。正则性取 $\dim R=\dim_\kappa\mathfrak m/\mathfrak m^2$ 的定义。标准先修为正则局部环的参数构成正则序列及 Cohen--Macaulay 性质、深度引理、有限自由消解的基本比较映射、Nakayama 引理、由有限模检测 Noetherian 环整体维数的判据（10.109.12），以及有限自由复形的子式--深度正合性判据（10.102.9；其完整作者证明已在本库 006 中附出）。正文实质完成 Koszul 复形向极小消解的比较和两个维数界，没有直接引用本等价定理。10.110.3 末尾的线性单射验证按作者原文留为可直接检验的步骤，未补写助手证明。全部支撑引理只计入本问题。

保留作者原语言、论证次序及数学内容；仅调整独立文档的排版和外部引用显示。原有章节号均指源书，不是本文件编号。网页评论不属于作者证明，未收录。
\subsection*{可修改的简短标注（非作者正文）}
$c$：局部环正则性的同调刻画

$a$：Koszul 复形；外幂；极小自由消解；剩余域；余切空间；Cohen--Macaulay 模；深度与子式理想；Nakayama 引理

\texttt{cross\_theory\_status = yes}。把余切空间的维数通过 Koszul 复形嵌入极小消解，再以深度控制消解长度，建立局部几何维数与同调维数的等同。位置：10.110.3 的比较映射和 10.110.5。

全部标注均为非穷尽、可修改初稿；未标注不表示未使用。
\subsection*{许可与修改记录}
原数学正文为 The Stacks Project Authors 的作品，按 GNU Free Documentation License 1.2 或后续版本复用；无不变章节、封面文字或封底文字。许可全文在 \texttt{sources/stacks/GFDL-1.2.txt}。本整理沿用原许可。2026-09-28：助手摘录、独立排版并加入来源、范围说明与初稿标注；不暗示原作者认可本次整理。
\end{document}
